\documentclass[11pt]{amsart}

\usepackage[margin=1.1in]{geometry}
\usepackage{amsmath,amssymb,amsthm,mathtools}
\usepackage{enumitem}
\usepackage[colorlinks=true,linkcolor=blue,urlcolor=blue]{hyperref}

\numberwithin{equation}{section}

\theoremstyle{definition}
\newtheorem{definition}{Definition}[section]
\newtheorem{example}[definition]{Example}
\newtheorem{exercise}[definition]{Exercise}
\theoremstyle{plain}
\newtheorem{proposition}[definition]{Proposition}
\newtheorem{lemma}[definition]{Lemma}
\theoremstyle{remark}
\newtheorem{remark}[definition]{Remark}

\newcommand{\R}{\mathbb{R}}
\newcommand{\E}{\mathbb{E}}
\newcommand{\Prob}{\mathbb{P}}
\newcommand{\cD}{\mathcal{D}}
\newcommand{\cH}{\mathcal{H}}
\newcommand{\cL}{\mathcal{L}}
\newcommand{\cP}{\mathcal{P}}
\newcommand{\cR}{\mathcal{R}}
\newcommand{\cS}{\mathcal{S}}
\newcommand{\cX}{\mathcal{X}}
\newcommand{\cY}{\mathcal{Y}}
\newcommand{\dd}{\,\mathrm{d}}
\newcommand{\one}{\mathbf{1}}
\newcommand{\norm}[1]{\left\lVert#1\right\rVert}
\newcommand{\ip}[2]{\left\langle#1,#2\right\rangle}
\DeclareMathOperator{\KL}{KL}
\DeclareMathOperator{\TV}{TV}
\DeclareMathOperator{\softmax}{softmax}
\DeclareMathOperator*{\argmin}{arg\,min}
\DeclareMathOperator*{\argmax}{arg\,max}

\title[How Do Large Language Models Work?]
{How Do Large Language Models Work?\\
A Mathematical Guide and Connections to Scientific and World Models}
\author{Fei Lu}
\date{Fall 2026. Disclaimer: the notes are written with the help from ChatGPT,
and the author takes full responsibility for any errors. If you find any
errors, please contact the author at
\href{mailto:feilu@math.jhu.edu}{feilu@math.jhu.edu}.}

\begin{document}

\begin{abstract}
An autoregressive large language model is a parameterized family of
next-token conditional distributions.  A tokenizer supplies the discrete
state space, a causally masked Transformer computes the context-dependent
logits, cross-entropy training fits the conditional law, and a decoding rule
turns that law into text.  Its analysis must separate the population target,
finite architecture, optimization algorithm, and deployment rule.  We give an exact
three-token example in which one attention head implements Bayesian
in-context learning, and then replace tokens by functions, trajectories, and
controlled states.  The resulting links to scientific foundation models and
world models expose additional problems of identifiability, discretization,
intervention, planning, and rollout stability.
\end{abstract}

\maketitle

\section{The mathematical object behind a language model}
\label{sec:language-law}

\paragraph{\textbf{Key question.}}
What object does an LLM learn, what computation evaluates it, and why can that
same computation appear to adapt to a prompt without changing its parameters?

The main chain in this note is
\begin{equation}\label{eq:main-chain}
\begin{gathered}
\text{raw text}
\longrightarrow \text{tokens}
\longrightarrow \text{causal hidden states}
\longrightarrow \text{next-token law}\\
\longrightarrow \text{cross-entropy training}
\longrightarrow \text{autoregressive generation}
\longrightarrow \text{prompt-conditioned action}.
\end{gathered}
\end{equation}
Each arrow introduces a different mathematical approximation.  In particular,
a theorem about the population log-loss minimizer is not a theorem about
finite-data generalization, optimization of a Transformer, or the quality of a
decoded answer.

\subsection{Tokenization and finite strings}

Let $\mathcal V$ be a finite vocabulary with $K=|\mathcal V|$ (the K in practice is often order $10^4$ to $10^5$).
It contains special symbols such as beginning-of-sequence and end-of-sequence tokens.  A
tokenizer is a deterministic map
\[
 \operatorname{Tok}:\{\text{strings}\}\longrightarrow
 \bigcup_{T\ge 1}\mathcal V^T.
\]
Modern tokenizers often use subword or byte-level units; byte-pair encoding is
one influential construction \cite{SennrichEtAl2016}.  The representation is
not mathematically neutral: the vocabulary determines sequence length, the
units on which likelihood is measured, and which patterns are local.

For clarity, first fix a maximum horizon $T$ and include an end-of-sequence
token.  Let $U_{1:T}=(U_1,\ldots,U_T)$ have population law $p_*$ on
$\mathcal V^T$.  Represent a shorter string by its first end token and repeat
that token thereafter; training may mask the repeated positions.  The
probability chain rule gives
\begin{equation}\label{eq:chain-rule}
 p_*(u_{1:T})
 =
 \prod_{t=1}^T p_*(u_t\mid u_{<t}),
 \qquad
 u_{<t}=(u_1,\ldots,u_{t-1}).
\end{equation}
Thus a joint law on strings can be represented by a collection of next-token
conditionals.

\begin{definition}[Autoregressive language model]
\label{def:autoregressive-lm}
An autoregressive language model is a parameterized family
\[
 p_\theta(\,\cdot\mid u_{<t})\in\cP(\mathcal V),
 \qquad \theta\in\Theta,
\]
for every admissible prefix $u_{<t}$.  It induces the joint law
\begin{equation}\label{eq:model-joint-law}
 p_\theta(u_{1:T})
 =\prod_{t=1}^T p_\theta(u_t\mid u_{<t}).
\end{equation}
\end{definition}

The learned object is therefore a conditional probability vector, not a
unique answer.  Sampling, greedy decoding, and search are later decisions that
use this vector in different ways.

\subsection{Why next-token log loss is the natural population target}

Let $q(\,\cdot\mid u_{<t})$ be any candidate conditional law that is positive
on the support of $p_*(\,\cdot\mid u_{<t})$.  Define its population sequence
log loss by
\begin{equation}\label{eq:population-log-loss}
 \cL(q)
 =
 \E_{p_*}\left[
   -\sum_{t=1}^T\log q(U_t\mid U_{<t})
 \right].
\end{equation}

\begin{proposition}[Conditional log loss identifies the token law]
\label{prop:log-loss-kl}
Assume that the terms below are finite.  Then
\begin{equation}\label{eq:conditional-kl-sum}
 \cL(q)-\cL(p_*)
 =
 \sum_{t=1}^T
 \E\!\left[
   \KL\!\left(
     p_*(\,\cdot\mid U_{<t})
     \,\middle\|\,
     q(\,\cdot\mid U_{<t})
   \right)
 \right]
 \ge 0.
\end{equation}
Equality holds if and only if the two conditional distributions agree
$p_*$-almost surely at every predicted position.
\end{proposition}

\begin{proof}
Condition on $U_{<t}$ and add and subtract
$-\log p_*(U_t\mid U_{<t})$.  The conditional expectation of the difference is
the displayed Kullback--Leibler divergence.  Sum over $t$ and use its
nonnegativity.
\end{proof}

This proposition explains the statistical target of pretraining in an
unrestricted model class.  It does not say that a finite Transformer contains
the true law, that a finite corpus determines it, or that a nonconvex optimizer
finds the best parameter.

The average token loss is often reported through perplexity,
\[
 \operatorname{PPL}(q)
 =\exp\!\left\{\frac{1}{T}\cL(q)\right\}.
\]
Perplexity depends on the tokenizer.  Moreover, the population distribution is
a distribution of text, not a distribution of truth.  Low next-token loss
alone does not imply factuality, calibrated uncertainty, valid reasoning, or
good decisions under a downstream loss.


\section{One forward pass through a decoder-only Transformer}
\label{sec:transformer-forward}

\paragraph{\textbf{Key question.}}
How does one fixed network map a variable token prefix to the conditional
probability vector in Definition~\ref{def:autoregressive-lm}?

Let the input contain $N$ tokens $u_1,\ldots,u_N$.  We reserve $N$ for the
number of tokens or particles, $L$ for the number of Transformer layers, and
$d$ for the hidden width.  Let $e_{u_i}\in\R^K$ be the one-hot vector of token
$u_i$.  The initial hidden state is
\begin{equation}\label{eq:input-embedding}
 x_i^0=E_{\rm tok}e_{u_i}+r_i+s_i\in\R^d,
 \qquad i=1,\ldots,N,
\end{equation}
where $E_{\rm tok}\in\R^{d\times K}$ is the token-embedding matrix, $r_i$ is
a positional representation, and $s_i$ may encode a role or modality.  Learned
absolute positions, sinusoidal encodings, and rotary representations are
different implementations of the same need: token order must enter the model.

\subsection{Causal self-attention}

At layer $\ell$ and attention head $h$, let
\[
 Q_{\ell h},K_{\ell h}\in\R^{d_k\times d},
 \qquad
 V_{\ell h}\in\R^{d_v\times d}.
\]
Given normalized states $\bar x_i^\ell$, define the query, key, and value
vectors
\[
 q_i^{\ell h}=Q_{\ell h}\bar x_i^\ell,
 \qquad
 k_j^{\ell h}=K_{\ell h}\bar x_j^\ell,
 \qquad
 v_j^{\ell h}=V_{\ell h}\bar x_j^\ell.
\]
The causal attention weights and head output are
\begin{equation}\label{eq:causal-attention}
\begin{aligned}
  z_i^{\ell h} =\sum_{j=1}^i a_{ij}^{\ell h}v_j^{\ell h}, \qquad 
  a_{ij}^{\ell h} =
 \frac{
  \exp\!\left(\ip{q_i^{\ell h}}{k_j^{\ell h}}/\sqrt{d_k}\right)
  \one_{\{j\le i\}}
 }{
  \displaystyle\sum_{m=1}^i
  \exp\!\left(\ip{q_i^{\ell h}}{k_m^{\ell h}}/\sqrt{d_k}\right)
 }.
\end{aligned}
\end{equation}
For each $i$, the weights are nonnegative and sum to one.  Thus a single head
output lies in the convex hull of its value vectors.  The output projection,
residual connection, and feed-forward layer need not preserve this convexity.

For $H$ heads, concatenate their outputs and apply
$W_O^\ell\in\R^{d\times Hd_v}$.  Let $d_{\rm ff}$ be the feed-forward width,
with
\[
 W_1^\ell\in\R^{d_{\rm ff}\times d},\quad
 b_1^\ell\in\R^{d_{\rm ff}},\quad
 W_2^\ell\in\R^{d\times d_{\rm ff}},\quad
 b_2^\ell\in\R^d.
\]
A simplified pre-normalized residual block is
\begin{equation}\label{eq:transformer-block}
\begin{aligned}
 \bar x_i^\ell
 &=\operatorname{LN}_\ell(x_i^\ell),\\
 \widetilde x_i^\ell
 &=x_i^\ell
   +W_O^\ell\operatorname{Concat}
      (z_i^{\ell1},\ldots,z_i^{\ell H}),\\
 x_i^{\ell+1}
 &=\widetilde x_i^\ell
   +W_{2}^\ell
     \sigma\!\left(
       W_{1}^\ell\operatorname{LN}'_\ell(\widetilde x_i^\ell)
       +b_{1}^\ell
     \right)
   +b_{2}^\ell.
\end{aligned}
\end{equation}
Layer normalization acts across the $d$ features of each token.  In its basic
form,
\[
 \operatorname{LN}(x)
 =g\odot
   \frac{x-\bar x\one_d}{\sqrt{d^{-1}\norm{x-\bar x\one_d}^2+\varepsilon}}
   +b,
 \qquad
 \bar x=d^{-1}\sum_{j=1}^d x_j,
\]
with learned $g,b\in\R^d$ and fixed $\varepsilon>0$.  Practical architectures
vary in their normalization, nonlinearities, biases, attention variants, and
parameter sharing, but the causal dependency is invariant across these
choices.

After $L$ layers, the logits at position $i$ are
\begin{equation}\label{eq:output-softmax}
 g_i=W_{\rm out}\operatorname{LN}_{\rm out}(x_i^L)+b_{\rm out}\in\R^K,
 \qquad
 p_\theta(v\mid u_{\le i})
 =\frac{e^{(g_i)_v}}{\sum_{w\in\mathcal V}e^{(g_i)_w}}, 
\end{equation}
where $W_{\rm out}\in\R^{K\times d}$ and $b_{\rm out}\in\R^K$.
The same parameter vector $\theta$ contains the embeddings, all attention and
feed-forward weights, normalizations, and the output map.

When matrix notation is convenient below, we stack token states as rows of
$X\in\R^{N\times d}$.  Part II uses the column-token convention
$X^\top\in\R^{d\times N}$; the two descriptions are transposes of one
another.

\subsection{Two exact structural properties}

\begin{proposition}[Causal dependence]
\label{prop:causal-dependence}
For the architecture in \eqref{eq:input-embedding}--\eqref{eq:transformer-block},
$x_i^\ell$ is a function only of $(u_1,\ldots,u_i)$ and the fixed marks at
positions $1,\ldots,i$.  In particular, the logits used to predict $U_{i+1}$
do not depend on future tokens.
\end{proposition}

\begin{proof}
The assertion holds for $\ell=0$ by \eqref{eq:input-embedding}.  Assume it
holds at layer $\ell$.  The causal mask in \eqref{eq:causal-attention} restricts
$z_i^{\ell h}$ to states $\bar x_j^\ell$ with $j\le i$, each of which depends
only on $u_{\le j}\subseteq u_{\le i}$.  The output projection, normalization,
residual addition, and feed-forward map act at position $i$.  They cannot
introduce dependence on a later token.  Induction over $\ell$ proves the claim.
\end{proof}

This property permits teacher-forced training: all positions can be processed
in parallel even though each prediction uses only its true past.

\begin{proposition}[Permutation equivariance without marks]
\label{prop:permutation-equivariance}
Remove positional and role marks and replace the causal mask by unmasked
self-attention.  If $P$ is an $N\times N$ permutation matrix acting on token
rows, then every Transformer block satisfies
\[
 \operatorname{Block}(PX)=P\operatorname{Block}(X).
\]
Thus an unmarked, unmasked stack cannot distinguish two sequences that differ
by a permutation through any permutation-invariant readout.  More precisely,
the stack carries no absolute order information: permuting the input rows only
permutes the output rows.
\end{proposition}

\begin{proof}
Let $S(X)=(XQ^\top)(XK^\top)^\top/\sqrt{d_k}$ be the score matrix for one
head, using row-token notation.  Then $S(PX)=PS(X)P^\top$.  Row-wise softmax
commutes with simultaneous row and column permutation, so
\[
 \softmax(S(PX))=P\softmax(S(X))P^\top.
\]
Since $(PX)V^\top=P(XV^\top)$, the attention output is permuted by $P$.
Concatenation,
row-wise normalization, pointwise feed-forward maps, and residual addition all
preserve the equivariance.
\end{proof}

Positions and causal masks are therefore not cosmetic.  They convert a
set-valued interaction into an ordered computation.

\subsection{The exact interacting-particle identity}

The link to Part II begins before any limit.  For one unmasked head, define the
empirical token measure
\[
 \mu_X^N=\frac1N\sum_{j=1}^N\delta_{x_j}.
\]
Then the attention output is exactly
\begin{equation}\label{eq:attention-empirical-measure}
 \operatorname{Att}_\theta(X)_i
 =
 \frac{\displaystyle
   \int e^{\ip{Qx_i}{Ky}/\sqrt{d_k}}Vy\,\mu_X^N(\dd y)
 }{\displaystyle
   \int e^{\ip{Qx_i}{Kz}/\sqrt{d_k}}\,\mu_X^N(\dd z)
 }.
\end{equation}
For language, an unmarked measure loses order.  Introduce the marked empirical
measure
\[
 \eta_X^N=\frac1N\sum_{j=1}^N\delta_{(j,x_j)}.
\]
The causal head at position $i$ has the exact marked-measure representation
\begin{equation}\label{eq:causal-attention-marked-measure}
 \operatorname{Att}^{\rm causal}_\theta(X)_i
 =
 \frac{\displaystyle
   \int \one_{\{m\le i\}}
   e^{\ip{Qx_i}{Ky}/\sqrt{d_k}}Vy\,
   \eta_X^N(\dd m,\dd y)
 }{\displaystyle
   \int \one_{\{m\le i\}}
   e^{\ip{Qx_i}{Ky}/\sqrt{d_k}}\,
   \eta_X^N(\dd m,\dd y)
 }.
\end{equation}
This is an exact finite-token identity.  Replacing the empirical measure by a
limiting law, sending the layer count to a continuous-depth limit, or taking a
high-dimensional limit are separate mathematical steps.


\section{How the parameters are learned}
\label{sec:training}

\paragraph{\textbf{Key question.}}
What does the training algorithm optimize, and which conclusions follow from
that objective alone?

Let the pretraining corpus be
\[
 \cD_{\rm pre}=\{u_{1:T_r}^{(r)}:r=1,\ldots,R\},
 \qquad
 S=\sum_{r=1}^R T_r.
\]
The empirical next-token loss is
\begin{equation}\label{eq:empirical-pretraining-loss}
 \widehat\cL_S(\theta)
 =
 -\frac1S
 \sum_{r=1}^R\sum_{t=1}^{T_r}
 \log p_\theta\!\left(
   u_t^{(r)}\mid u_{<t}^{(r)}
 \right).
\end{equation}
The ground-truth prefix is supplied at every training position; this is
\emph{teacher forcing}.  Causal masking makes all the losses in a document
computable in one parallel forward pass.

\subsection{The local softmax calculation}

For logits $g\in\R^K$ and target class $y$, the single-token loss is
\begin{equation}\label{eq:softmax-loss}
 \ell(g;y)
 =-\log\frac{e^{g_y}}{\sum_{j=1}^K e^{g_j}}
 =-g_y+\log\sum_{j=1}^K e^{g_j}.
\end{equation}

\begin{proposition}[Gradient and curvature of the output loss]
\label{prop:softmax-gradient}
Let $p=\softmax(g)$ and let $e_y$ be the $y$th coordinate vector.  Then
\begin{equation}\label{eq:softmax-gradient}
 \nabla_g\ell(g;y)=p-e_y,
 \qquad
 \nabla_g^2\ell(g;y)=\operatorname{diag}(p)-pp^\top\succeq0.
\end{equation}
\end{proposition}

\begin{proof}
Differentiating \eqref{eq:softmax-loss} gives
$\partial_{g_j}\ell=p_j-\one_{\{j=y\}}$.  A second derivative gives the
matrix $\operatorname{diag}(p)-pp^\top$.  For $a\in\R^K$,
\[
 a^\top(\operatorname{diag}(p)-pp^\top)a
 =\sum_jp_ja_j^2-\left(\sum_jp_ja_j\right)^2
 =\operatorname{Var}_{J\sim p}(a_J)\ge0.
\]
\end{proof}

The loss is convex in the output logits but not in the full parameter vector
$\theta$.  Backpropagation applies the chain rule through every operation in
Section~\ref{sec:transformer-forward}.  A schematic minibatch update is
\begin{equation}\label{eq:sgd-update}
 \theta_{k+1}
 =\theta_k-\gamma_k\widehat G_k,
 \qquad
 \widehat G_k
 =\frac1{|B_k|}\sum_{b\in B_k}\nabla_\theta\ell_b(\theta_k),
\end{equation}
with learning rate $\gamma_k$.  If $B_k$ is sampled uniformly from a fixed
finite collection of token losses, then $\widehat G_k$ is an unbiased gradient
of that finite empirical loss.  Adam and related methods replace the scalar
step by adaptive moment-based preconditioning.  None of these observations
gives a global convergence theorem for a production-scale Transformer.

Empirical scaling studies found regular power-law behavior of language-model
loss over substantial ranges of model size, data, and compute
\cite{KaplanEtAl2020,HoffmannEtAl2022}.  These are empirical relations within
specified regimes, not universal learning-theory laws and not explanations of
semantic competence.

\subsection{Why a chat model is not merely a base language model}

Next-token pretraining imitates the distribution of the corpus.  Instruction
tuning instead minimizes a supervised loss on prompt--response pairs.  A later
preference stage may compare candidate responses and optimize a learned reward;
InstructGPT is an influential example \cite{OuyangEtAl2022}.

One useful idealization starts from a base response law $p_0(\,\cdot\mid c)$
and a bounded reward $r(c,y)$.  For a fixed prompt $c$, consider
\begin{equation}\label{eq:kl-regularized-preference}
 \sup_{\pi(\cdot\mid c)}
 \left\{
   \E_{Y\sim\pi(\cdot\mid c)}[r(c,Y)]
   -\beta\KL\!\left(
      \pi(\,\cdot\mid c)\|p_0(\,\cdot\mid c)
    \right)
 \right\},
 \qquad \beta>0.
\end{equation}

\begin{proposition}[An idealized preference tilt]
\label{prop:gibbs-tilt}
If
$Z(c)=\E_{p_0(\cdot\mid c)}[e^{r(c,Y)/\beta}]<\infty$, the optimizer of
\eqref{eq:kl-regularized-preference} is
\begin{equation}\label{eq:gibbs-tilt}
 \pi^*(y\mid c)
 =\frac{p_0(y\mid c)e^{r(c,y)/\beta}}{Z(c)}.
\end{equation}
\end{proposition}

\begin{proof}
Define $\pi^*$ by \eqref{eq:gibbs-tilt}.  For any $\pi\ll p_0$,
\[
 \E_\pi[r]-\beta\KL(\pi\|p_0)
 =\beta\log Z(c)-\beta\KL(\pi\|\pi^*).
\]
The last term is minimized uniquely at $\pi=\pi^*$.
\end{proof}

This identity explains how a reward can tilt a pretrained distribution while a
KL penalty discourages excessive departure from it.  Actual preference
pipelines must estimate the reward and approximately optimize a parameterized
policy, so the proposition is a benchmark rather than a description of every
post-training algorithm.


\section{From logits to generated text}
\label{sec:generation}

\paragraph{\textbf{Key question.}}
How does a conditional probability vector become one finite output string?

Given a prompt $c=u_{1:N}$, full-softmax ancestral sampling repeats
\begin{equation}\label{eq:ancestral-step}
 U_{N+t}\sim p_\theta(\,\cdot\mid U_{1:N+t-1})
\end{equation}
until an end token or a prescribed horizon is reached.

\begin{proposition}[Ancestral sampling realizes the product law]
\label{prop:ancestral-law}
Conditional on the prompt $c$, the procedure in
\eqref{eq:ancestral-step} generates
\begin{equation}\label{eq:conditional-product-law}
 \Prob_\theta(U_{N+1:N+T}=y_{1:T}\mid c)
 =\prod_{t=1}^T
   p_\theta(y_t\mid c,y_{<t}).
\end{equation}
\end{proposition}

\begin{proof}
Apply the conditional probability chain rule to the sequential sampling
procedure.  At step $t$, its transition probability is precisely the $t$th
factor on the right.
\end{proof}

If $g_v$ is the model logit for token $v$, temperature decoding uses
\begin{equation}\label{eq:temperature-law}
 p_{\theta,\tau_{\rm dec}}(v\mid c)
 =\frac{e^{g_v/\tau_{\rm dec}}}
        {\sum_{w\in\mathcal V}e^{g_w/\tau_{\rm dec}}},
 \qquad \tau_{\rm dec}>0.
\end{equation}
At $\tau_{\rm dec}=1$ this is the trained softmax.  As
$\tau_{\rm dec}\downarrow0$, it concentrates on the largest logit when the
maximizer is unique.  Top-$k$ and nucleus sampling first truncate the support
and then renormalize.  Greedy decoding takes the locally most probable token,
while beam search approximately searches several prefixes.  These rules change
the output law; greedy decoding does not in general find the globally most
probable sequence.

Training conditions on true prefixes, whereas free-running generation
conditions on previously generated tokens.  A small one-step average loss can
therefore coexist with rare but consequential errors along a long generated
trajectory.  The same issue will reappear for scientific rollouts and world
models.

At inference, previously computed keys and values can be cached.  Generating a
new token then attends to the cached prefix rather than recomputing every old
hidden state.  This reduces redundant computation, but exact dense attention
still compares the new token with all preceding tokens.


\section{A complete miniature model of in-context learning}
\label{sec:toy-icl}

\paragraph{\textbf{Key question.}}
Can a fixed attention mechanism exactly change its prediction using only a
short prompt?

The following example separates representability from optimization.  It is
small enough that the Bayes rule and the attention weights can both be written
exactly.

\begin{example}[A noisy two-task language]
\label{ex:noisy-task}
Let a latent task $W\in\{-1,+1\}$ be uniform.  Conditional on $W$, observe
$Y_1,Y_2,Y_3\in\{-1,+1\}$ independently with
\begin{equation}\label{eq:noisy-task-law}
 \Prob(Y_i=W\mid W)=\frac34,
 \qquad
 \Prob(Y_i=-W\mid W)=\frac14.
\end{equation}
A pretraining episode is serialized as $(Y_1,Y_2,Y_3,W)$.  The final
next-token problem is to predict the task token $W$ from the three-token
prompt.
\end{example}

Let $S_3=Y_1+Y_2+Y_3$.  The posterior log odds are
\begin{equation}\label{eq:toy-posterior-odds}
 \log\frac{\Prob(W=+1\mid Y_{1:3})}
          {\Prob(W=-1\mid Y_{1:3})}
 =S_3\log 3.
\end{equation}
Consequently,
\begin{equation}\label{eq:toy-posterior-values}
\begin{array}{c|cccc}
S_3&3&1&-1&-3\\ \hline
\Prob(W=+1\mid Y_{1:3})
&27/28&3/4&1/4&1/28.
\end{array}
\end{equation}

\begin{proposition}[One attention head realizes the Bayes rule]
\label{prop:toy-attention-bayes}
There is a width-two, one-head causal attention model with fixed weights whose
next-token softmax after $(Y_1,Y_2,Y_3)$ equals the posterior in
\eqref{eq:toy-posterior-values} for every prompt.
\end{proposition}

\begin{proof}
For this exact toy construction, omit layer normalization and the
feed-forward sublayer; retain one attention sublayer, its residual connection,
and the final linear readout.
Embed token $Y_i$ as $x_i=(Y_i,0)\in\R^2$.  Set queries and keys to zero, so at
position $3$ the causal attention weights on positions $1,2,3$ are all $1/3$.
Choose the value map
\[
 V(y,0)=(0,y).
\]
The head output is $(0,m)$ with $m=S_3/3$.  After the residual connection, the
hidden state is $(Y_3,m)$; choose the output map to ignore the first coordinate
and set
\[
 g_+(m)=\frac32(\log3)m,
 \qquad
 g_-(m)=-\frac32(\log3)m.
\]
Then
\[
 g_+(m)-g_-(m)=3(\log3)m=S_3\log3,
\]
which equals the posterior log odds in \eqref{eq:toy-posterior-odds}.  The
two-class softmax therefore gives the exact Bayes posterior.
\end{proof}

Greedy decoding implements majority vote, whereas sampling retains posterior
uncertainty.  The network weights do not change when the three prompt tokens
change.  This is literal in-context adaptation.  It is only a representability
result: it does not say that gradient descent must discover these weights or
that an open-ended LLM performs exact Bayesian inference.

\subsection{The general latent-task benchmark}

Let $W\sim\Pi$ index a task or rule.  A fresh context contains $n$
demonstrations
\[
 C_n=((Z_1,Y_1),\ldots,(Z_n,Y_n)),
\]
and a query $Z_*$.  Suppose that $Z_*$ is fixed by the experiment and
$Y_*\mathrel{\perp\!\!\!\perp}C_n\mid(W,Z_*)$.  For every measurable
$A\subseteq\cY$, the posterior predictive law is
\begin{equation}\label{eq:general-posterior-predictive}
 \Prob(Y_*\in A\mid C_n,Z_*)
 =\int
   \Prob_w(Y_*\in A\mid Z_*)\,
   \Pi(\dd w\mid C_n).
\end{equation}
If the query also contains information about $W$, the posterior must condition
on $(C_n,Z_*)$.

Write $M$ for the number of pretraining tasks or episodes and $n$ for the
number of demonstrations from a fresh task.  A trained Transformer defines an
action
\begin{equation}\label{eq:icl-action-risk}
 \widehat a_*
 =\mathsf T_{\widehat\theta_M}(C_n;Z_*),
 \qquad
 \cR_{M,n}
 =\E\!\left[
   \ell(\widehat a_*,Y_*)
 \right].
\end{equation}
The expectation includes the pretraining sample and the fresh task, context,
query, and response.  Increasing $M$ estimates shared structure; increasing
$n$ supplies information about the realized task.  They address different
uncertainties.

\begin{proposition}[More nested context reduces Bayes log risk]
\label{prop:context-entropy}
Assume $Y_*$ is discrete and integrable log losses are finite.  If
$\sigma(C_n,Z_*)\subseteq\sigma(C_{n+1},Z_*)$, then
\begin{equation}\label{eq:entropy-monotonicity}
 H(Y_*\mid C_{n+1},Z_*)
 \le H(Y_*\mid C_n,Z_*).
\end{equation}
Equivalently, the expected Bayes log loss cannot increase when an additional
context observation is revealed.
\end{proposition}

\begin{proof}
The difference is the conditional mutual information
\[
 H(Y_*\mid C_n,Z_*)-H(Y_*\mid C_{n+1},Z_*)
 =I(Y_*;C_{n+1}\mid C_n,Z_*)\ge0,
\]
where $C_{n+1}$ denotes the augmented nested context.
\end{proof}

This is an oracle statement.  A learned finite network may fail to use the
additional information and can even perform worse with a poorly designed or
out-of-distribution prompt.  Xie et al.~\cite{XieEtAl2022} prove an
implicit-Bayes mechanism for a particular mixture-of-Hidden-Markov-models
pretraining distribution; GPT-3 supplies the canonical large-scale empirical
observation of fixed-weight few-shot prompting \cite{BrownEtAl2020}.


\section{From language models to scientific foundation models}
\label{sec:scientific-fm}

\paragraph{\textbf{Key question.}}
What changes when tokens represent functions, fields, equations, and solution
operators rather than words?

Let $\cX$ and $\cY$ be separable Hilbert spaces.  A scientific task $W=w$ may
determine a solution operator
\begin{equation}\label{eq:scientific-operator}
 \cS_w:\cX\longrightarrow\cY,
\end{equation}
for an ODE, PDE, inverse problem, or controlled dynamical system.  A fresh
context can contain noisy probes
\begin{equation}\label{eq:scientific-context}
 C_n=\{(f_i,\cS_wf_i+\xi_i)\}_{i=1}^n,
 \qquad \xi_i\in\cY,
\end{equation}
along with symbolic equations, coefficients, geometry, boundary conditions,
sensor locations, or natural-language metadata.  Given a query $f_*$, a
pretrained model returns
\begin{equation}\label{eq:scientific-action}
 \widehat u_*
 =\mathsf T_{\widehat\theta_M}(C_n;f_*).
\end{equation}
The output may instead be an unknown coefficient, a governing equation, a
future trajectory, or a conditional law.  These are different estimands and
require different losses.

The analogy with an LLM is precise at the level of conditional learning:
\begin{equation}\label{eq:llm-science-comparison}
\begin{aligned}
\text{LLM:}\quad
&p_\theta(u_{t+1}\mid u_{\le t}),\\
\text{scientific model:}\quad
&\mathsf P_\theta(\dd u\mid C_n,f_*),
\qquad u=\cS_Wf_*.
\end{aligned}
\end{equation}
Pretraining amortizes computation across a distribution of documents or
scientific tasks; context identifies the part of that shared structure relevant
to the query.

The scientific setting adds structure absent from generic token prediction:
\begin{enumerate}[label=(\arabic*)]
\item $\cS_w$ acts between function spaces, so input and output resolutions
$(p,q)$ introduce approximation and discretization error.
\item The context may not identify $w$ or even $\cS_wf_*$; inverse stability and
the probe information operator determine what is observable.
\item Conservation laws, symmetries, boundary conditions, semigroup structure,
and well-posedness can restrict the admissible operator class.
\item A small pointwise or one-step loss need not control a solver, inverse
reconstruction, or long rollout.
\end{enumerate}

\subsection{Four distinct meanings of adaptation}

It is important not to place all scientific foundation models under the label
``in-context learning.''

\paragraph{\textbf{Data-prompted operator adaptation.}}
ICON, introduced by Yang, Liu, Meng, and Osher
\cite{YangEtAl2023ICON}, uses condition--response function pairs to specify a
new operator and answers a query without updating the weights.  This is the
closest analogue of GPT-style ICL.  The demonstrated operators belong to
controlled synthetic families, and the work does not prove a universal
operator-learning rate.

\paragraph{\textbf{Symbolically specified PDE prediction.}}
PDEformer-1, developed by Ye et al. with Bin Dong
\cite{YeEtAl2024PDEformer}, represents a known one-dimensional PDE as a graph,
combines symbolic and numerical inputs with a graph Transformer, and queries an
implicit solution representation.  It demonstrates zero-shot prediction in
declared families and data-efficient fine-tuning for new equations.  The
symbolic PDE is part of the input; this is not context-only identification from
probe pairs.

\paragraph{\textbf{Multimodal equation and trajectory prediction.}}
PROSE \cite{LiuEtAl2024PROSE} jointly uses numerical trajectories and symbolic
ODEs to forecast solutions and generate governing expressions across a
synthetic equation family.  Its zero-shot scope is relative to that generator.

\paragraph{\textbf{Pretraining followed by fine-tuning.}}
Poseidon \cite{HerdeEtAl2024Poseidon} pretrains an operator Transformer on
fluid-dynamics solution operators and transfers it by fine-tuning to downstream
PDE tasks.  Its use of multiple temporal pairs exploits semigroup structure;
transfer after fine-tuning is not the same as fixed-weight ICL.

\subsection{Why rollout stability is a separate theorem}

Suppose a scientific model learns a one-step map $\widehat S$ intended to
approximate a true evolution map $S$.  Repeated application is a deterministic
world model.

\begin{proposition}[Error accumulation for an iterated operator]
\label{prop:deterministic-rollout}
Let $D\subseteq\cX$ be invariant under $S$ and $\widehat S$.  Assume
\begin{equation}\label{eq:rollout-assumptions}
 \norm{Sx-Sy}\le L_S\norm{x-y},
 \qquad
 \sup_{x\in D}\norm{\widehat Sx-Sx}\le\delta.
\end{equation}
Then for every $x\in D$ and integer $k\ge1$,
\begin{equation}\label{eq:rollout-bound}
 \norm{\widehat S^kx-S^kx}
 \le
 \delta\sum_{j=0}^{k-1}L_S^j
 =
 \begin{cases}
 \delta(L_S^k-1)/(L_S-1),&L_S\ne1,\\
 k\delta,&L_S=1.
 \end{cases}
\end{equation}
\end{proposition}

\begin{proof}
Let $e_k=\norm{\widehat S^kx-S^kx}$.  Invariance and
\eqref{eq:rollout-assumptions} give
\[
\begin{aligned}
 e_{k+1}
 &\le
 \norm{\widehat S(\widehat S^kx)-S(\widehat S^kx)}
 +\norm{S(\widehat S^kx)-S(S^kx)}\\
 &\le\delta+L_Se_k.
\end{aligned}
\]
Since $e_0=0$, iteration yields \eqref{eq:rollout-bound}.
\end{proof}

For $L_S<1$, the bound is uniform in time; for $L_S=1$, it grows linearly;
for $L_S>1$, the worst-case estimate grows exponentially.  Directly learning a
long-time solution operator and iterating a one-step model are therefore
different approximation problems.


\section{World models: prediction under intervention and planning}
\label{sec:world-models}

\paragraph{\textbf{Key question.}}
When does a sequence predictor become a model that an agent can use to imagine
and choose future actions?

Consider a partially observed controlled process with latent state $S_t$,
action $A_t$, observation $O_t$, and reward $R_t$:
\begin{equation}\label{eq:pomdp-law}
\begin{aligned}
 S_{t+1}&\sim K_*(\,\cdot\mid S_t,A_t),\\
 O_t&\sim G_*(\,\cdot\mid S_t),\\
 R_t&\sim \rho_*(\,\cdot\mid S_t,A_t).
\end{aligned}
\end{equation}
The observable history is
$H_t=(O_1,A_1,R_1,\ldots,A_{t-1},R_{t-1},O_t)$.  A world model learns a
predictive object that can be rolled forward under proposed actions.  It may
model observations directly,
\begin{equation}\label{eq:observable-world-model}
 p_\theta(O_{t+1},R_t\mid H_t,A_t),
\end{equation}
or introduce a learned latent state $Z_t$ with transition, observation, and
reward models
\begin{equation}\label{eq:latent-world-model}
 p_\theta(Z_{t+1}\mid Z_t,A_t),
 \qquad
 p_\theta(O_t\mid Z_t),
 \qquad
 p_\theta(R_t\mid Z_t,A_t).
\end{equation}
An inference model $q_\phi(Z_t\mid H_t)$ estimates the current latent state
from history.

\begin{definition}[World model for a fixed interface]
\label{def:world-model}
Write $Y_{t+1}=(R_t,O_{t+1})$.  A world model is a family of
action-conditioned predictive kernels
\[
 \widehat P_\theta(\dd y_{t+1}\mid h_t,a_t)
\]
that can be recursively composed with a policy: a predicted outcome is
appended to the history and used to predict the next outcome.  The definition
alone does not assert that a latent coordinate is physical, that the learned
kernel identifies causal effects outside the support of the training actions,
or that it transfers to every environment.
\end{definition}

The distinction between a one-step fit and a complete imagined trajectory has
an exact information-theoretic form.  Let $\nu$ be the initial observation law
and define the true observable predictive kernel by
\[
 P_*^{\rm obs}(\dd y_{t+1}\mid h_t,a_t)
 =\operatorname{Law}(Y_{t+1}\in\dd y_{t+1}\mid H_t=h_t,A_t=a_t).
\]
Let $\pi(\dd a_t\mid h_t)$ be any possibly history-dependent policy.  The
true and learned laws of a horizon-$J$ trajectory are
\begin{align}
 \mathbf P_{*,J}^{\pi}(\dd h_{J+1})
 &=\nu(\dd o_1)\prod_{t=1}^J
   \pi(\dd a_t\mid h_t)P_*^{\rm obs}(\dd y_{t+1}\mid h_t,a_t),
 \label{eq:true-trajectory-law}\\
 \widehat{\mathbf P}_{\theta,J}^{\pi}(\dd h_{J+1})
 &=\nu(\dd o_1)\prod_{t=1}^J
   \pi(\dd a_t\mid h_t)\widehat P_\theta
   (\dd y_{t+1}\mid h_t,a_t).
 \label{eq:learned-trajectory-law}
\end{align}

\begin{proposition}[Trajectory KL chain rule]
\label{prop:trajectory-kl}
Suppose that the true one-step kernel is absolutely continuous with respect to
the learned kernel at every history--action pair reached under
$\mathbf P_{*,J}^{\pi}$.  Then, with values in $[0,\infty]$,
\begin{equation}\label{eq:trajectory-kl}
 \KL\!\left(
   \mathbf P_{*,J}^{\pi}
   \,\middle\|\,
   \widehat{\mathbf P}_{\theta,J}^{\pi}
 \right)
 =\sum_{t=1}^J
  \E_{\mathbf P_{*,J}^{\pi}}
  \left[
   \KL\!\left(
    P_*^{\rm obs}(\,\cdot\mid H_t,A_t)
    \,\middle\|\,
    \widehat P_\theta(\,\cdot\mid H_t,A_t)
   \right)
  \right].
\end{equation}
\end{proposition}

\begin{proof}
In the likelihood ratio of
\eqref{eq:true-trajectory-law} with respect to
\eqref{eq:learned-trajectory-law}, the common initial law and policy factors
cancel.  Hence the log likelihood ratio is the sum of the $J$ conditional log
likelihood ratios.  Taking expectation under the true trajectory law and then
conditioning on $(H_t,A_t)$ gives \eqref{eq:trajectory-kl}.
\end{proof}

Thus trajectory KL is a sum of one-step errors, but the errors are evaluated
under the state--action occupancy of the policy being assessed.  Small
training loss under one behavior policy does not control a new planning policy
without coverage, uniform approximation, or another change-of-measure
argument.

\subsection{Latent-state training}

Consider a trajectory of length $T$.  For observations $O_{1:T}$ and
action--reward pairs $(A_t,R_t)_{t=1}^{T-1}$, one possible factorization is
\begin{equation}\label{eq:world-generative-factorization}
\begin{aligned}
&p_\theta(z_{1:T},o_{1:T},r_{1:T-1}\mid a_{1:T-1})\\
&\qquad={}
 p_\theta(z_1)
 \prod_{t=1}^T p_\theta(o_t\mid z_t)
 \prod_{t=1}^{T-1}
 p_\theta(r_t\mid z_t,a_t)
 p_\theta(z_{t+1}\mid z_t,a_t).
\end{aligned}
\end{equation}
A sequential variational posterior of the form
\begin{equation}\label{eq:world-variational-posterior}
 q_\phi(z_{1:T}\mid H_T)
 =q_\phi(z_1\mid H_1)
  \prod_{t=1}^{T-1}
  q_\phi(z_{t+1}\mid z_t,H_{t+1})
\end{equation}
yields the evidence lower bound
\begin{equation}\label{eq:world-elbo}
\begin{aligned}
 \log p_\theta(o_{1:T},r_{1:T-1}\mid a_{1:T-1})
 \ge{}&
 \E_{q_\phi}\!\left[
   \sum_{t=1}^T\log p_\theta(o_t\mid Z_t)
   +\sum_{t=1}^{T-1}\log p_\theta(r_t\mid Z_t,a_t)
 \right]\\
 &-\KL(q_\phi(Z_1\mid H_1)\|p_\theta(Z_1))\\
 &-\sum_{t=1}^{T-1}
   \E_{q_\phi(Z_t\mid H_t)}\!\left[
     \KL\!\left(
       q_\phi(Z_{t+1}\mid Z_t,H_{t+1})
       \|p_\theta(Z_{t+1}\mid Z_t,a_t)
     \right)
   \right].
\end{aligned}
\end{equation}
This is one template, not a definition.  PlaNet uses a recurrent state-space
model for latent planning \cite{HafnerEtAl2019Planet}; Dreamer learns behavior
from imagined latent rollouts \cite{HafnerEtAl2020Dreamer}.  MuZero instead
learns latent dynamics, rewards, values, and policies sufficient for tree
search without requiring pixel reconstruction \cite{SchrittwieserEtAl2020}.
Thus a useful world-model state need not be a faithful physical coordinate; it
may retain only decision-relevant information.

\subsection{Planning uses the model outside a one-step loss}

For a planning horizon $H_{\rm plan}$, a model-based agent evaluates proposed
actions by
\begin{equation}\label{eq:model-return}
 \widehat J_{H_{\rm plan}}(a_{t:t+H_{\rm plan}-1};H_t)
 =\E_\theta\!\left[
   \sum_{k=0}^{H_{\rm plan}-1}\gamma^kR_{t+k}
   \,\middle|\,
   H_t,a_{t:t+H_{\rm plan}-1}
 \right].
\end{equation}
It may optimize the action sequence online, train a policy inside imagined
trajectories, or combine a learned model with tree search.  The relevant error
is therefore not merely the average one-step prediction loss under the training
distribution.

The following elementary simulation bound makes the point precise.  Define
$\TV(\mu,\nu)=\sup_A|\mu(A)-\nu(A)|$.

\begin{proposition}[A discounted simulation bound]
\label{prop:simulation-lemma}
Consider two fully observed discounted MDPs with the same reward
$0\le r(s,a)\le R_{\max}$, discount $0<\gamma<1$, and transition kernels
$P$ and $\widehat P$.  If
\begin{equation}\label{eq:uniform-tv-error}
 \sup_{s,a}\TV(P(\,\cdot\mid s,a),\widehat P(\,\cdot\mid s,a))
 \le\varepsilon,
\end{equation}
then for every stationary policy $\pi$,
\begin{equation}\label{eq:simulation-bound}
 \norm{V_P^\pi-V_{\widehat P}^\pi}_\infty
 \le
 \frac{\gamma R_{\max}}{(1-\gamma)^2}\,\varepsilon.
\end{equation}
\end{proposition}

\begin{proof}
Let $P_\pi$ and $\widehat P_\pi$ be the policy-averaged kernels and let
$r_\pi$ be the common expected reward.  The Bellman equations give
\[
 V_P^\pi-V_{\widehat P}^\pi
 =\gamma P_\pi(V_P^\pi-V_{\widehat P}^\pi)
  +\gamma(P_\pi-\widehat P_\pi)V_{\widehat P}^\pi.
\]
Since $0\le V_{\widehat P}^\pi\le R_{\max}/(1-\gamma)$ and mixing over actions
preserves the total-variation bound,
\[
 \norm{(P_\pi-\widehat P_\pi)V_{\widehat P}^\pi}_\infty
 \le\frac{R_{\max}}{1-\gamma}\varepsilon.
\]
The layer-cake representation gives the constant one, rather than two,
because $V_{\widehat P}^\pi$ is nonnegative and bounded and because we use the
convention $\TV(\mu,\nu)=\sup_A|\mu(A)-\nu(A)|$.
Taking sup norms, moving the term
$\gamma\norm{V_P^\pi-V_{\widehat P}^\pi}_\infty$ to the left, and dividing by
$1-\gamma$ proves the result.
\end{proof}

The uniform hypothesis in \eqref{eq:uniform-tv-error} is strong.  A model
trained on trajectories from one behavior policy may be inaccurate under the
state--action distribution produced by a new planner.  This is the world-model
version of distribution shift and model exploitation.

\subsection{From recurrent world models to foundation world models}

Ha and Schmidhuber \cite{HaSchmidhuber2018} learned compressed visual states
and recurrent stochastic dynamics, then trained a small controller using the
learned model; their experiments established the modern terminology in a
limited set of game environments.  Later systems separated several meanings
of a world model:
\begin{enumerate}[label=(\arabic*)]
\item a reconstructive latent simulator, as in PlaNet and Dreamer;
\item a decision-sufficient latent model, as in MuZero;
\item a large generative interactive environment pretrained from broad video
data, as in Genie \cite{BruceEtAl2024Genie}.
\end{enumerate}
Genie combines a video tokenizer, autoregressive dynamics, and a latent-action
model learned from unlabeled gameplay videos.  It is a useful example of a
``foundation world model,'' but its reported scope is still far from a general
physical simulator: the original system has short memory, slow generation,
and can produce unrealistic futures.

An LLM is not automatically a world model.  Its token law is trained on
observational text, while a world model for control must predict consequences
under interventions $A_t=a$.  Text correlations need not identify such causal
effects; a hidden state need not correspond to a physical state; and fluent
rollouts need not be dynamically calibrated.  A multimodal LLM can nevertheless
be part of a world-model system when it is grounded in observations, actions,
and feedback and when its predictive state is tested under intervention.


\section{One mathematical spine for the three model classes}
\label{sec:unification}

The common object is a conditional transition or response law:
\begin{equation}\label{eq:three-laws}
\begin{aligned}
\text{language model:}\quad
&p_\theta(u_{t+1}\mid u_{\le t}),\\
\text{scientific foundation model:}\quad
&\mathsf P_\theta(\dd u\mid C_n,f_*),
\quad u=\cS_Wf_*,\\
\text{world model:}\quad
&\widehat P_\theta(\dd y_{t+1}\mid h_t,a_t).
\end{aligned}
\end{equation}
All three amortize prediction over a distribution of related situations.  They
differ in what is represented, what information is supplied at deployment,
and which downstream stability statement is required.

\paragraph{\textbf{Representation.}}
An LLM discretizes text into tokens.  A scientific model discretizes functions,
meshes, graphs, sensors, and equations.  A world model compresses observation
histories into predictive latent states.  Representation error is therefore
task dependent and can remove information before learning begins.

\paragraph{\textbf{Conditional target.}}
The target may be a next-token law, a point solution, an operator value, a
posterior distribution, a reward, or a decision-sufficient state.  Log loss,
squared prediction loss, operator norm, conservation error, and control value
are not interchangeable.

\paragraph{\textbf{Deployment.}}
Adaptation may occur through a context with fixed weights, through explicit
symbolic conditioning, through fine-tuning, or through online interaction.
These mechanisms use different information and require different risk
definitions.

\paragraph{\textbf{Stability.}}
For language, generated tokens become future inputs.  For PDE and world models,
predicted states become future states used by solvers or planners.  In every
case, a local regression theorem must be paired with a stability or simulation
argument before it becomes an end-to-end guarantee.

\subsection{Connection to the rest of the course}

\paragraph{\textbf{Part II: attention as an interacting particle system.}}
Equation~\eqref{eq:attention-empirical-measure} is the exact finite system.
Part II studies its mean-field, continuous-depth, stability, clustering, and
inverse-interaction questions while keeping positions and masks as marks.

\paragraph{\textbf{Part III: generative modeling as learned dynamics.}}
An LLM generates a discrete path by ancestral sampling.  Score, diffusion, and
flow models generate continuous states by learned ODEs or SDEs.  Part III asks
how regression error in a learned field controls the resulting law and how
conditioning produces posterior samples for inverse problems.

\paragraph{\textbf{Part IV: ICL, inverse problems, and operator learning.}}
The latent-task model in \eqref{eq:general-posterior-predictive} becomes a
statistical experiment.  Part IV studies context observability, Bayes and
empirical-Bayes risk, task diversity, the separate roles of $M$ and $n$, and
the passage from finite vectors to operators at resolutions $(p,q)$.

The resulting learning-theory program is not a universal additive error
formula.  For each model one must specify:
\begin{enumerate}[label=(\arabic*)]
\item the task or environment distribution and the deployment distribution;
\item the estimand and downstream loss;
\item the information supplied by pretraining, context, symbols, sensors, and
actions;
\item approximation and optimization errors of the architecture;
\item discretization, rollout, inverse, or planning stability converting local
prediction into final risk.
\end{enumerate}


\section{Exercises with solutions}
\label{sec:exercises}

\begin{exercise}[Chain rule with an end token]
Let $\mathcal V$ contain an end token $\mathtt{EOS}$.  Suppose a random finite
string has length at most $T$ and, after its first end token, remains at
$\mathtt{EOS}$.  Show that its law is determined by the next-token conditionals
and derive \eqref{eq:chain-rule}.  Explain why the end token is needed for a
normalized law on variable-length strings.
\end{exercise}

\begin{proof}[Solution]
For every $u_{1:T}$ with positive probability, repeated conditioning gives
\[
 \Prob(U_{1:T}=u_{1:T})
 =\Prob(U_1=u_1)
  \prod_{t=2}^T
  \Prob(U_t=u_t\mid U_{<t}=u_{<t}).
\]
This is \eqref{eq:chain-rule}, with the empty prefix used for $t=1$.  Conversely,
suppose every conditional is normalized and impose the absorbing condition
\[
 q(\mathtt{EOS}\mid u_{<t})=1
 \quad\text{whenever $u_{<t}$ already contains $\mathtt{EOS}$.}
\]
Summing the product recursively over $u_T,u_{T-1},\ldots,u_1$ gives one, and
the absorbing condition restricts its support to padded finite strings.  The
first $\mathtt{EOS}$ identifies the finite stopping time.  Without an end
convention, a law on finite strings would also require separate probabilities
for termination at every prefix.
\end{proof}

\begin{exercise}[Causality and permutation]
Prove directly for one attention head that unmasked attention without position
marks is permutation equivariant.  Then give a two-token example showing that
adding distinct positional vectors can distinguish $(u_1,u_2)$ from
$(u_2,u_1)$.
\end{exercise}

\begin{proof}[Solution]
Let $A(X)=\softmax(S(X))XV^\top$, with tokens stored as rows and softmax applied
row-wise.  For a permutation matrix $P$,
$S(PX)=PS(X)P^\top$ and hence
\[
 A(PX)
 =P\softmax(S(X))P^\top P(XV^\top)
 =PA(X).
\]
For a concrete scalar example, take token embeddings $e_u=1$, $e_v=-1$ and
positional marks $r_1=1$, $r_2=-1$.  The marked representations of $(u,v)$
and $(v,u)$ are respectively
\[
 X_{uv}=(2,-2)^\top,
 \qquad
 X_{vu}=(0,0)^\top.
\]
If $P$ exchanges the two rows, then $X_{vu}\ne PX_{uv}$.  Thus even the
identity pointwise block followed by a fixed-position readout can distinguish
the two orders.  The point is that positions stay fixed when token contents
are exchanged; they do not travel with the token permutation.
\end{proof}

\begin{exercise}[Softmax curvature]
Derive both identities in \eqref{eq:softmax-gradient}.  Show that the Hessian
has $\one_K$ in its null space and interpret this invariance.
\end{exercise}

\begin{proof}[Solution]
Differentiation gives
$\partial_{g_j}\ell=p_j-\one_{\{j=y\}}$ and
$\partial_{g_k}p_j=p_j(\one_{\{j=k\}}-p_k)$, which yields the Hessian.
Moreover,
\[
 (\operatorname{diag}(p)-pp^\top)\one_K
 =p-p(p^\top\one_K)=p-p=0.
\]
Adding the same constant to every logit leaves softmax unchanged, so the loss
is flat in the direction $\one_K$.
\end{proof}

\begin{exercise}[The exact three-token learner]
For Example~\ref{ex:noisy-task}, derive the posterior odds
\eqref{eq:toy-posterior-odds}, compute the four probabilities in
\eqref{eq:toy-posterior-values}, and verify that the attention construction in
Proposition~\ref{prop:toy-attention-bayes} gives the same probabilities.
\end{exercise}

\begin{proof}[Solution]
For one observation $y\in\{-1,+1\}$,
\[
 \frac{\Prob(Y_i=y\mid W=+1)}{\Prob(Y_i=y\mid W=-1)}
 =3^y.
\]
Conditional independence and equal prior odds give likelihood ratio
$3^{Y_1+Y_2+Y_3}=3^{S_3}$.  Therefore
\[
 \Prob(W=+1\mid Y_{1:3})
 =\frac{3^{S_3}}{1+3^{S_3}},
\]
which gives $27/28,3/4,1/4,1/28$ at $S_3=3,1,-1,-3$.  The attention head
returns $m=S_3/3$, and its logit difference is
$3(\log3)m=S_3\log3$.  The two-class softmax converts this logit difference to
$3^{S_3}/(1+3^{S_3})$, as required.
\end{proof}

\begin{exercise}[Rollout regimes]
Under the assumptions of Proposition~\ref{prop:deterministic-rollout}, simplify
the bound for $L_S<1$, $L_S=1$, and $L_S>1$.  What changes if the one-step error
at time $j$ is $\delta_j$ rather than a common $\delta$?
\end{exercise}

\begin{proof}[Solution]
If $L_S<1$, then
\[
 e_k\le\delta\frac{1-L_S^k}{1-L_S}
 \le\frac{\delta}{1-L_S}.
\]
For $L_S=1$, $e_k\le k\delta$.  For $L_S>1$,
$e_k\le\delta(L_S^k-1)/(L_S-1)$, which can grow exponentially.  If the error
at step $j$ is $\delta_j$, the recursion becomes
$e_{j+1}\le L_Se_j+\delta_j$, and induction yields
\[
 e_k\le\sum_{j=0}^{k-1}L_S^{k-1-j}\delta_j.
\]
Thus early errors receive the largest amplification when $L_S>1$.
\end{proof}


\section{A short reading path}

The following ten papers give a focused path through the three model classes.
For each one, identify the conditional target, the information supplied at
deployment, whether the weights change, and the downstream loss.
\begin{enumerate}[label=\textbf{R\arabic*.},leftmargin=*]
\item Vaswani et al.~\cite{VaswaniEtAl2017}: the Transformer architecture.
\item Brown et al.~\cite{BrownEtAl2020}: large-scale autoregressive ICL.
\item Xie et al.~\cite{XieEtAl2022}: a restricted implicit-Bayes mechanism.
\item Yang et al.~\cite{YangEtAl2023ICON}: fixed-weight in-context operator
learning.
\item Ye et al.~\cite{YeEtAl2024PDEformer}: symbolic--numerical PDE
pretraining.
\item Herde et al.~\cite{HerdeEtAl2024Poseidon}: pretrained solution
operators and fine-tuning.
\item Ha and Schmidhuber~\cite{HaSchmidhuber2018}: compressed recurrent world
models.
\item Hafner et al.~\cite{HafnerEtAl2019Planet}: latent dynamics for planning.
\item Schrittwieser et al.~\cite{SchrittwieserEtAl2020}: decision-sufficient
latent dynamics in MuZero.
\item Bruce et al.~\cite{BruceEtAl2024Genie}: a video-pretrained interactive
world model.
\end{enumerate}


\begin{thebibliography}{99}

\bibitem{BrownEtAl2020}
T.~B.~Brown et al.,
\emph{Language models are few-shot learners},
Advances in Neural Information Processing Systems \textbf{33} (2020),
1877--1901.
\href{https://proceedings.neurips.cc/paper/2020/hash/1457c0d6bfcb4967418bfb8ac142f64a-Abstract.html}
{NeurIPS proceedings}.

\bibitem{BruceEtAl2024Genie}
J.~Bruce, M.~Dennis, A.~Edwards, J.~Parker-Holder, Y.~Shi, E.~Hughes,
M.~Lai, A.~Mavalankar, R.~Steigerwald, C.~Apps, et al.,
\emph{Genie: Generative interactive environments},
Proceedings of the 41st International Conference on Machine Learning,
PMLR \textbf{235} (2024), 4603--4623.
\href{https://proceedings.mlr.press/v235/bruce24a.html}{PMLR version}.

\bibitem{HaSchmidhuber2018}
D.~Ha and J.~Schmidhuber,
\emph{Recurrent world models facilitate policy evolution},
Advances in Neural Information Processing Systems \textbf{31} (2018).
\href{https://papers.nips.cc/paper_files/paper/2018/hash/2de5d16682c3c35007e4e92982f1a2ba-Abstract.html}
{NeurIPS proceedings}.

\bibitem{HafnerEtAl2019Planet}
D.~Hafner, T.~Lillicrap, I.~Fischer, R.~Villegas, D.~Ha, H.~Lee, and
J.~Davidson,
\emph{Learning latent dynamics for planning from pixels},
Proceedings of the 36th International Conference on Machine Learning,
PMLR \textbf{97} (2019), 2555--2565.
\href{https://proceedings.mlr.press/v97/hafner19a.html}{PMLR version}.

\bibitem{HafnerEtAl2020Dreamer}
D.~Hafner, T.~Lillicrap, J.~Ba, and M.~Norouzi,
\emph{Dream to control: Learning behaviors by latent imagination},
International Conference on Learning Representations, 2020.
\href{https://openreview.net/forum?id=S1lOTC4tDS}{OpenReview version}.

\bibitem{HerdeEtAl2024Poseidon}
M.~Herde, B.~Raoni\'c, T.~Rohner, R.~K\"appeli, R.~Molinaro,
E.~de B\'ezenac, and S.~Mishra,
\emph{Poseidon: Efficient foundation models for PDEs},
Advances in Neural Information Processing Systems \textbf{37} (2024).
\href{https://proceedings.neurips.cc/paper_files/paper/2024/hash/84e1b1ec17bb11c57234e96433022a9a-Abstract-Conference.html}
{NeurIPS proceedings}.

\bibitem{HoffmannEtAl2022}
J.~Hoffmann et al.,
\emph{Training compute-optimal large language models},
Advances in Neural Information Processing Systems \textbf{35} (2022),
30016--30030.
\href{https://proceedings.neurips.cc/paper_files/paper/2022/hash/c1e2faff6f588870935f114ebe04a3e5-Abstract-Conference.html}
{NeurIPS proceedings}.

\bibitem{KaplanEtAl2020}
J.~Kaplan, S.~McCandlish, T.~Henighan, T.~B.~Brown, B.~Chess, R.~Child,
S.~Gray, A.~Radford, J.~Wu, and D.~Amodei,
\emph{Scaling laws for neural language models},
\href{https://arxiv.org/abs/2001.08361}{arXiv:2001.08361}, 2020.

\bibitem{LiuEtAl2024PROSE}
Y.~Liu, Z.~Zhang, and H.~Schaeffer,
\emph{PROSE: Predicting multiple operators and symbolic expressions using
multimodal Transformers},
Neural Networks \textbf{180} (2024), 106707.
\href{https://doi.org/10.1016/j.neunet.2024.106707}
{doi:10.1016/j.neunet.2024.106707}.

\bibitem{OuyangEtAl2022}
L.~Ouyang et al.,
\emph{Training language models to follow instructions with human feedback},
Advances in Neural Information Processing Systems \textbf{35} (2022),
27730--27744.
\href{https://arxiv.org/abs/2203.02155}{arXiv:2203.02155}.

\bibitem{SchrittwieserEtAl2020}
J.~Schrittwieser et al.,
\emph{Mastering Atari, Go, chess and shogi by planning with a learned model},
Nature \textbf{588} (2020), 604--609.
\href{https://doi.org/10.1038/s41586-020-03051-4}
{doi:10.1038/s41586-020-03051-4}.

\bibitem{SennrichEtAl2016}
R.~Sennrich, B.~Haddow, and A.~Birch,
\emph{Neural machine translation of rare words with subword units},
Proceedings of the 54th Annual Meeting of the Association for Computational
Linguistics (2016), 1715--1725.
\href{https://aclanthology.org/P16-1162/}{ACL Anthology}.

\bibitem{VaswaniEtAl2017}
A.~Vaswani, N.~Shazeer, N.~Parmar, J.~Uszkoreit, L.~Jones, A.~N.~Gomez,
\L.~Kaiser, and I.~Polosukhin,
\emph{Attention is all you need},
Advances in Neural Information Processing Systems \textbf{30} (2017),
5998--6008.
\href{https://proceedings.neurips.cc/paper_files/paper/2017/hash/3f5ee243547dee91fbd053c1c4a845aa-Abstract.html}
{NeurIPS proceedings}.

\bibitem{XieEtAl2022}
S.~M.~Xie, A.~Raghunathan, P.~Liang, and T.~Ma,
\emph{An explanation of in-context learning as implicit Bayesian inference},
International Conference on Learning Representations, 2022.
\href{https://openreview.net/forum?id=RdJVFCHjUMI}{OpenReview version}.

\bibitem{YangEtAl2023ICON}
L.~Yang, S.~Liu, T.~Meng, and S.~J.~Osher,
\emph{In-context operator learning with data prompts for differential equation
problems},
Proc. Natl. Acad. Sci. USA \textbf{120} (2023), no.~39, e2310142120.
\href{https://doi.org/10.1073/pnas.2310142120}
{doi:10.1073/pnas.2310142120}.

\bibitem{YeEtAl2024PDEformer}
Z.~Ye, X.~Huang, L.~Chen, Z.~Liu, B.~Wu, H.~Liu, Z.~Wang, and B.~Dong,
\emph{PDEformer-1: A foundation model for one-dimensional partial differential
equations},
\href{https://arxiv.org/abs/2407.06664}{arXiv:2407.06664}, 2024.

\end{thebibliography}

\end{document}
